\section{Causal Attribution Framework}
\label{sec:causal}

The foundational paper's scorpion detection
(Proposition~2 of \citet{lasser2026gfm}) identifies agents whose actions are
\emph{correlated} with $\volL$-contraction. It explicitly acknowledges that
``disentangling the agent's contribution from simultaneous causes requires a
causal-attribution procedure that the framework does not provide.'' This
section provides that procedure.

\subsection{Structural Causal Model for Capability Dynamics}

\begin{definition}[Capability Dynamics SCM]
\label{def:scm}
A \emph{structural causal model for capability dynamics} is a tuple
    $\SCM = (U, V, F, P(U))$ where:
\begin{itemize}
\item $U$: exogenous variables representing agent dispositions, environmental
    conditions, and stochastic factors not determined by the model.
\item $V = \{G_t, \pi_t^{(1)}, \ldots, \pi_t^{(n)}, \Delta \volL_t\}$:
    endogenous variables representing the capability poset state, each
    agent's action at time $t$, and the resulting $\volL$-change.
\item $F$: structural equations mapping parent variables to children in the
    causal directed acyclic graph (DAG). The DAG encodes which agents'
    actions causally influence which capability changes.
\item $P(U)$: prior distribution over exogenous variables.
\end{itemize}
When the variance-reduction results of this paper are invoked
    (Proposition~\ref{prop:causal_dominates}(c) and downstream), we
    additionally assume \textbf{exogenous orthogonality}: the
    exogenous noise components entering different agents' structural
    equations, and the residual $h$, are mutually independent
    conditional on $G_t$. That is,
    if $U_j \subseteq U$ are the exogenous variables appearing in
    $a_j$'s structural equation, $U_k \subseteq U$ those appearing
    in $a_k$'s, and $U_h \subseteq U$ those appearing in the
    residual $h(G_t, U)$
    (Equation~\eqref{eq:additive_scm}), then
    $U_j \perp\!\!\!\perp U_k \mid G_t$ for $j \neq k$ and
    $U_j \perp\!\!\!\perp U_h \mid G_t$ for all $j$.
    This rules out shared exogenous shocks that would induce
    covariance between agents' contributions (or between an agent's
    contribution and the environment residual) even under a correct
    causal DAG.

The GFM actor's world model $\WorldModel$
    (Definition~1 of \citet{lasser2026horizon}) is an instance of $\SCM$
    restricted to the actor's beliefs about causal structure.
\end{definition}

The SCM extends $\WorldModel$ by making the causal structure explicit. In the
foundational paper, $\WorldModel$ is an opaque prediction model: given
actions, it predicts $\volL$-changes. The SCM decomposes predictions into
causal pathways, enabling counterfactual queries: ``what would have happened
if agent $a_j$ had acted differently?''

\paragraph{Temporal scope.} Definition~\ref{def:scm} specifies a
\emph{single-step} SCM: the endogenous variables
$\{G_t, \pi_t^{(1)}, \ldots, \pi_t^{(n)}, \Delta \volL_t\}$ describe
one time slice, and the structural equations $F$ map actions at step
$t$ to the $\volL$-change at step $t$. The do-intervention
$\mathrm{do}(\pi_j = \cdot)$ operates within this single-step
structure: ``other variables respond according to $F$'' means other
endogenous variables in the same time slice are recomputed under the
intervened structural equations, not that the system evolves forward
in time. Multi-step effects (an action at $t$ that causes a cascade
at $t+1$) are captured by instantiating a new single-step SCM at
each subsequent time step with an updated poset state $G_{t+1}$;
the causal attribution $\CausalAttr(a_j, t)$
(Definition~\ref{def:causal_attribution}) is the single-step effect.

\subsection{Causal Contraction Attribution}

\begin{definition}[Causal Contraction Attribution]
\label{def:causal_attribution}
For an observed $\volL$-change $\Delta \volL(G_t)$ at time $t$, the
    \emph{causal contribution} of agent $a_j$ is:
\begin{equation}
\label{eq:causal_attribution}
\CausalAttr(a_j, t) = \Delta \volL\bigl(G_t \mid \mathrm{do}(\pi_j =
    \pi_j^{\mathrm{obs}})\bigr)
    - \Delta \volL\bigl(G_t \mid \mathrm{do}(\pi_j = \mathrm{skip})\bigr)
\end{equation}
where $\mathrm{do}(\pi_j = \cdot)$ denotes the do-calculus intervention
    \citep{pearl2009causality}: setting $a_j$'s action to a specific value
    while allowing all other variables to respond according to $F$. The
    first term is the $\volL$-change under $a_j$'s observed action; the
    second is the counterfactual $\volL$-change had $a_j$ done nothing.
\end{definition}

The causal attribution $\CausalAttr(a_j, t)$ isolates $a_j$'s contribution
from confounders: simultaneous actions by other agents, environmental
changes, and background dynamics that the correlational signal in
Proposition~2 of \citet{lasser2026gfm} cannot separate. When $a_j$'s
actions cause contraction, $\CausalAttr(a_j, t) < 0$. When they cause
expansion, $\CausalAttr(a_j, t) > 0$. When $a_j$'s actions have no causal
effect (the contraction would have occurred regardless), $\CausalAttr(a_j,
t) = 0$, and the correlational detection would have produced a false
positive.

\paragraph{Additivity.} The one-at-a-time $\CausalAttr(a_j, t)$ of
Equation~\eqref{eq:causal_attribution} sums to the total
$\volL$-change only when the SCM's structural equation for
$\Delta \volL$ is \emph{additive in agent actions}: there exist
functions $g_j$ and a residual $h$ such that
\begin{equation}
\label{eq:additive_scm}
\Delta \volL(G_t)
    = \sum_j g_j\bigl(\pi_j^{\mathrm{obs}}, G_t, U\bigr)
    + h(G_t, U)
\end{equation}
where $h$ does not depend on any $\pi_j$. Under this additive-SCM
assumption, the one-at-a-time intervention
$\mathrm{do}(\pi_j = \pi_j^{\mathrm{obs}})$ vs.\
$\mathrm{do}(\pi_j = \mathrm{skip})$ isolates $g_j$ exactly, and
\begin{equation}
\label{eq:causal_decomposition}
\Delta \volL(G_t) = \sum_j \CausalAttr(a_j, t) + h(G_t, U).
\end{equation}
The additive-SCM assumption is stronger than (and is not implied by)
conditional independence of agent actions given $G_t$: two
independently drawn actions can still interact multiplicatively in a
non-additive structural equation such as
$\Delta \volL = \pi_1 \cdot \pi_2$, in which case the one-at-a-time
$\CausalAttr$ does not sum to the total. When interaction effects
are present and a non-additive SCM is required, an exact decomposition
needs a Shapley-style attribution rule that averages $\CausalAttr$
over all orderings of agent inclusion \citep{shapley1953value}; the
one-at-a-time form is then an approximation whose error is the
deviation from additivity. For scorpion detection
(Section~\ref{sec:scorpion}), the one-at-a-time form is sufficient
because the detection signal depends only on the sign and magnitude
of the diagonal contribution $g_j$, not on the exact decomposition
of interaction effects.

\subsection{Causal Detection Dominates Correlational Detection}

\begin{proposition}[Convergence Advantage of Causal Attribution]
\label{prop:causal_dominates}
Let $a_j$ be an agent with expected causal contraction $\mu_j =
    \mathbb{E}[\CausalAttr(a_j, t)] < 0$ (a scorpion). Assume the
    observation sequences for both detectors are i.i.d.\ within a
    stationary strategy epoch, or weakly dependent with summable
    autocovariance. Under the SCM $\SCM$, the following
    \emph{information-theoretic} sample-complexity comparison holds
    (independent of the specific sequential test used; see
    Section~\ref{sec:scorpion} for the SPRT-specific operational
    rates):
\begin{enumerate}
\item[\textbf{(a)}] The correlational detector (Proposition~2(a) of
    \citet{lasser2026gfm}) has sample complexity
    $\Omega(\sigma^2 / \mu_j^2)$ for identifying $a_j$ as a
    net-contractor, under the
    idealized assumption that $\mu_j$ appears as an
    isolable drift in the trust-weighted aggregate $\Deltahat$ (see
    paragraph below on the correlational estimator). Here $\sigma^2$
    is the total variance of the signal the correlational detector
    can construct, including confounders.
\item[\textbf{(b)}] The causal detector using $\CausalAttr(a_j, t)$
    has sample complexity
    $\Omega(\sigma_{\mathrm{do}}^2 / \mu_j^2)$, where
    $\sigma_{\mathrm{do}}^2$ is the variance of the causal attribution
    signal after conditioning on the do-intervention. The ratio
    $\sigma_{\mathrm{do}}^2 / \sigma^2 \leq 1$ (under the conditions
    of part~(c)) is the structural advantage factor.
\item[\textbf{(c)}] Under \textbf{causal sufficiency} (the SCM
    $\SCM$ captures all common causes of $a_j$'s actions and the
    confounder signals), the \textbf{additive-SCM} assumption
    (Equation~\eqref{eq:additive_scm}), and \textbf{exogenous
    orthogonality} (Definition~\ref{def:scm}), we have
    $\mathrm{Cov}(\CausalAttr(a_j), \mathrm{confounders} \mid G_t) = 0$
    and therefore
    \[
        \sigma^2 - \sigma_{\mathrm{do}}^2
        = \mathrm{Var}(\text{confounders} \mid G_t) \geq 0,
    \]
    so $\sigma_{\mathrm{do}}^2 \leq \sigma^2$, with strict inequality
    whenever confounders contribute nonzero variance. Without causal
    sufficiency (or without the additive-SCM structure), the
    variance identity gives
    \[
        \sigma^2 - \sigma_{\mathrm{do}}^2
        = \mathrm{Var}(\text{confounders} \mid G_t)
        + 2\,\mathrm{Cov}(\CausalAttr(a_j), \mathrm{confounders} \mid G_t),
    \]
    whose sign depends on the covariance term. The inequality
    $\sigma_{\mathrm{do}}^2 \leq \sigma^2$ therefore requires the
    full triple (causal sufficiency, additive SCM, exogenous
    orthogonality), or the weaker condition
    $\mathrm{Cov}(\CausalAttr(a_j), \mathrm{confounders} \mid G_t)
        \geq -\tfrac{1}{2} \mathrm{Var}(\text{confounders} \mid G_t)$.
    Under a general (non-sufficient) SCM with adversarial confounder
    correlation, causal attribution can in principle achieve worse
    convergence than the correlational detector; characterizing
    populations in which this occurs is outside the scope of this
    paper.
\end{enumerate}
\end{proposition}

\begin{proof}
(a) The correlational detector's observable signal is the
trust-weighted aggregate
$\Deltahat(\pi_t) = \sum_{k \in \mathcal{O}} T_k \cdot R_k(\pi_t)$
(Equation~10 of \citet{lasser2026gfm}), where $R_k \in \{-1, 0, +1\}$
is a ternary directional vote. Let
$\Delta \volL_{\mathrm{total}}(t)$ denote the \emph{true} continuous
$\volL$-change at step $t$ that $\Deltahat$ estimates. The
correlational detector's sample complexity for detecting $a_j$'s
mean effect $\mu_j$ is governed by the variance of the continuous
quantity it must recover, not by the variance of the ternary
estimator itself: a sign-correct ternary estimator inherits the
variance of the underlying continuous signal it tracks, plus
discretization noise. Formally, $\sigma^2$ in the bound is the
variance of any estimator of the total $\volL$-change that the
correlational detector can construct from its observation channel.
The $\Omega(\sigma^2 / \mu_j^2)$ lower bound is the
information-theoretic sample complexity for detecting a mean shift
$\mu_j$ in noise of variance $\sigma^2$.

(b) The causal detector evaluates $\CausalAttr(a_j, t)$ directly by
comparing the observed $\volL$-change against the counterfactual under
$\mathrm{do}(\pi_j = \mathrm{skip})$. The variance of this signal is
$\sigma_{\mathrm{do}}^2 = \mathrm{Var}(\CausalAttr(a_j, t))$, which
excludes the confounding variance.

(c) The three assumptions do distinct work; we separate them.

\emph{Step 1: Additive SCM $\Rightarrow$ $\CausalAttr(a_j) = g_j$
exactly.} Under the additive-SCM assumption
(Equation~\eqref{eq:additive_scm}), the continuous $\volL$-change
decomposes as
$\Delta \volL_{\mathrm{total}}(t) = \CausalAttr(a_j, t) +
    \mathrm{confounders}(t)$,
where the confounder term aggregates all other agents' contributions
$\{g_k\}_{k \neq j}$ and the residual $h(G_t, U)$, and
$\CausalAttr(a_j, t) = g_j(\pi_j^{\mathrm{obs}}, G_t, U) -
g_j(\mathrm{skip}, G_t, U)$ is exactly $a_j$'s separable
contribution. Without additivity, $\CausalAttr(a_j)$ captures only
the diagonal term and misses interaction effects. By the variance
identity for sums:
\[
\sigma^2 = \mathrm{Var}(\CausalAttr(a_j))
    + \mathrm{Var}(\mathrm{confounders})
    + 2\,\mathrm{Cov}(\CausalAttr(a_j), \mathrm{confounders}).
\]

\emph{Step 2: Causal sufficiency $\Rightarrow$ confounders are
SCM-known.} Causal sufficiency ensures the DAG $\SCM$ captures all
common causes of $a_j$'s actions and the confounder signals. Every
variable contributing to the covariance term appears explicitly in
the structural equations, so the covariance can in principle be
computed from the DAG. Without causal sufficiency, latent common
causes contribute unmodeled covariance that may have either sign.

\emph{Step 3: Exogenous orthogonality $\Rightarrow$
$\mathrm{Cov}$ term vanishes.} Under exogenous orthogonality
(Definition~\ref{def:scm}), the noise $U_j$ driving $g_j$ is
independent of both the noise $\{U_k\}_{k \neq j}$ driving
$\{g_k\}_{k \neq j}$ and the noise $U_h$ driving the residual
$h$, conditional on $G_t$. Since the confounders term is
$\sum_{k \neq j} g_k + h$, each cross-covariance
$\mathrm{Cov}(g_j, g_k \mid G_t) = 0$ and
$\mathrm{Cov}(g_j, h \mid G_t) = 0$. This gives
$\mathrm{Cov}(\CausalAttr(a_j), \mathrm{confounders} \mid G_t) = 0$,
and therefore
$\sigma^2 - \sigma_{\mathrm{do}}^2 = \mathrm{Var}(\mathrm{confounders})
\geq 0$
as stated. Without exogenous orthogonality, the covariance term is
retained and the inequality
$\sigma_{\mathrm{do}}^2 \leq \sigma^2$ holds only when
$\mathrm{Var}(\mathrm{confounders}) +
2\,\mathrm{Cov}(\CausalAttr(a_j), \mathrm{confounders}) \geq 0$; the
inequality is strict whenever confounders have positive variance and
the covariance term does not dominate.
\end{proof}

\paragraph{The correlational detector as an estimator.} Part (a) of
Proposition~\ref{prop:causal_dominates} treats the correlational
detector as if $\mu_j$ appeared directly as a component of the mean
of the aggregate signal $\Deltahat$. This is an idealization: the
actual mean of $\Deltahat$ is the sum of all agents' mean effects
plus background dynamics, not $\mu_j$ alone. A correlational
estimator that recovered $\mu_j$ from $\Deltahat$ would need an
additional auxiliary procedure --- for example, regressing $\Deltahat$
against an indicator for $a_j$'s action times, which incurs its own
bias when confounders are correlated with $a_j$'s action schedule.
The $O(\sigma^2 / \mu_j^2)$ bound therefore describes the best-case
correlational detector: one that has access to $a_j$'s action schedule
and uses it as a conditioning variable. A weaker correlational
detector that only observes $\Deltahat$ without conditioning has a
strictly worse (and generally incomputable) rate. The causal detector
does not need this auxiliary procedure because
$\CausalAttr(a_j, t)$ already isolates $a_j$'s contribution by
construction.

\paragraph{Practical implication.} In populations with many simultaneously
acting agents, confounding variance can dominate the correlational signal.
A single scorpion among 100 simultaneously acting agents contributes
$\sim$1\% of the signal variance; the correlational detector needs
$\sim$100$\times$ more observations than the causal detector to achieve the
same confidence. The causal detector's advantage grows linearly with the
number of simultaneously acting confounders.

\subsection{Causal Attribution as One Channel Among Several}
\label{sec:multi_channel}

Proposition~\ref{prop:causal_dominates} assumes that the SCM $\SCM$ is
accurate enough to support do-calculus attribution, and that the
structural assumptions of part (c) (causal sufficiency and
additive-SCM structure) hold. In practice the SCM is learned from
observational data, partially specified, and adversarially exposed:
agents can strategically align their actions to mask causal
structure, and new agents or delegation chains can emerge during the
trajectory. A framework that depended on the SCM being perfectly
learned would be brittle. Instead, the framework is robust because
causal attribution is \emph{one channel} among several in a
multi-channel attribution system, and the system's overall
attribution accuracy degrades gracefully as any single channel's
reliability degrades.

Multi-channel attribution is the standard methodology in adversarial
intrusion analysis and intelligence tradecraft. The Diamond Model of
intrusion analysis \citep{caltagirone2013diamond} organizes attribution
around four analytic vertices --- adversary, capability,
infrastructure, victim --- linked by relationships, with confidence
coming from cross-vertex correlation rather than from any single
vertex. The \emph{Pyramid of Pain} \citep{bianco2013pyramid} ranks
observable signals by how costly they are for an adversary to change:
hash values, IP addresses, and domain names are cheap to rotate;
tools and tactics, techniques, and procedures (TTPs) are expensive to
retrain. The MITRE ATT\&CK framework \citep{strom2018attack} provides
a shared capability taxonomy that grounds TTP-level attribution in a
common vocabulary. The \emph{Analysis of Competing Hypotheses}
\citep{heuer1999psychology} provides a structured method for
comparing attribution hypotheses by diagnostic evidence, weighting
signals by their power to discriminate among candidates rather than
by their individual strength. In all four traditions, the robustness
of an attribution decision is driven by signals the adversary cannot
cheaply fake, and no single signal is load-bearing.

The GFM sequence already supplies several attribution channels,
distributed across papers:
\begin{itemize}
\item \textbf{Behavioral prediction residual.} The trust model of
    \citet{lasser2026gfm} computes $r_k(t) = R_k(t) - \hat{R}_k(t
    \mid M_k(G))$, the divergence between an agent's observed
    behavior and its predictive model. An agent whose actions
    consistently match its model concentrates on high $T_k$;
    systematic divergence drives $T_k$ down. This is a moderately
    expensive channel to fake: a scorpion must maintain behavioral
    consistency with a fictitious model to evade it.
\item \textbf{Correlational co-occurrence.} The trust-weighted
    $\Deltahat$ signal of \citet{lasser2026gfm} (Proposition~2(a))
    detects agents whose actions temporally co-occur with
    $\volL$-contraction. This is a cheap channel to fake --- an
    adversary can introduce uncorrelated dummy actions --- and is
    strictly dominated by the causal channel under the conditions of
    Proposition~\ref{prop:causal_dominates}(c).
\item \textbf{Capability gating.} Papers 2 and 3 of the GFM sequence
    \citep{lasser2026poset, lasser2026horizon} together provide a
    structural constraint: an agent cannot cause the contraction of
    a capability it does not have upstream access to, regardless of
    any other signal. This is the \emph{means-analysis} from the
    intrusion-analysis tradition and sits at the top of the
    robustness hierarchy: it reflects the structure of the poset
    itself rather than behavior an adversary can stage.%
    \footnote{The gate's strength is bounded by the fidelity of the
    maintained capability model: hidden delegation, latent substitute
    capabilities, or a misspecified poset can defeat it.}
    An attribution decision that conflicts with capability gating
    should be overridden by it, not aggregated with it.
\item \textbf{Causal contraction attribution} (this paper,
    Definition~\ref{def:causal_attribution}). Moderately expensive to
    fake, since the SCM demands counterfactual consistency, and
    strictly informative under the conditions of
    Proposition~\ref{prop:causal_dominates}(c). Vulnerable to SCM
    misspecification.
\item \textbf{SPRT scorpion detection} (this paper,
    Section~\ref{sec:scorpion}). Moderately expensive to fake under
    stationary epochs; vulnerable to fast adaptation as characterized
    by the evasion bandwidth bound.
\end{itemize}

Table~\ref{tab:channels} summarizes the channels by evasion cost and
required assumptions.

\begin{table}[h]
\centering
\small
\begin{tabular}{@{}lllp{3.8cm}@{}}
\toprule
\textbf{Channel} & \textbf{Cost to fake} & \textbf{Paper} & \textbf{Key assumptions} \\
\midrule
Capability gating  & Very high & 2, 3 & Trusted, complete capability graph \\
Causal attribution & Moderate  & 4    & Causal sufficiency, additive SCM, exog.\ orthogonality \\
Behavioral residual & Moderate & 1    & Predictive-model access \\
SPRT detection     & Moderate  & 4    & Gaussian, stationary epochs \\
Correlational co-occ. & Low   & 1    & (None beyond observation) \\
\bottomrule
\end{tabular}
\caption{Attribution channels ordered by adversarial evasion cost (highest first), following the Pyramid of Pain principle.}
\label{tab:channels}
\end{table}

An attribution decision should aggregate evidence across channels,
weighted by each channel's calibrated reliability, rather than rely
on the output of any single channel. (The aggregation is Bayesian in
spirit --- each channel contributes evidence toward or against a
suspect --- but this paper does not formalize the joint hypothesis
space or channel-dependence model; the full fusion rule is an open
direction identified below.) The multi-channel framing has three
architectural consequences:

\begin{enumerate}
\item \textbf{Capability gating is a hard filter, not a prior.}
    An agent that cannot structurally reach capability $c_i$ cannot
    have caused its contraction, and no amount of causal or
    correlational evidence should override that. The correct
    aggregation is a gate-then-aggregate architecture: apply the
    capability gate first to produce the candidate suspect set, then
    aggregate the remaining channels over the candidates.
\item \textbf{No single channel must be perfect.} Each channel has
    its own reliability, and the aggregate's accuracy depends on
    the joint reliability profile rather than on the
    maximum-reliability channel. This is the property that makes
    multi-channel attribution robust in the threat-intelligence
    setting and that makes the GFM actor robust to imperfect
    structure learning (Section~\ref{sec:discussion}).
\item \textbf{Channel weights should be calibrated from
    measurable proxies}, not fixed at design time. The causal
    channel's weight should scale with an SCM-confidence estimate
    derived from structure-learning diagnostics; the behavioral
    channel's weight should scale with a population-level
    reliability measure derived from the per-agent trust scores
    $T_k$ of \citet{lasser2026gfm} (distinct from the actor's
    self-trust $T_s$, which governs learning rate rather than
    channel reliability); the correlational channel's weight
    should scale with the confounding variance estimated from the
    population structure. The calibration procedure is outside the
    scope of this paper and is identified as an open direction in
    Section~\ref{sec:discussion}.
\end{enumerate}

Paper 4's contribution to this framework is the causal channel: a
principled, do-calculus-based attribution signal with a convergence
advantage over the correlational baseline under stated structural
conditions. It is not the attribution system, and the paper does not
claim that the causal channel is load-bearing. The overall
attribution robustness comes from the multi-channel structure, and
from capability gating in particular as the structural floor.
